\documentclass[../../main.tex]{subfiles}

\begin{document}

\chapter[Symbol Table dan Type Checking]{Symbol Table, Scope Management,\\ dan Type Checking}
\label{chap:symbol-table}

\begin{outcomebox}
\textbf{Sub-CPMK:} \subcpmk{1.6}\\[3pt]
Setelah menyelesaikan bab ini, mahasiswa mampu:
\begin{enumerate}[label=\arabic*.,leftmargin=*]
  \item menjelaskan struktur dan operasi tabel simbol;
  \item mengimplementasikan tabel simbol dengan scope bersarang;
  \item menjelaskan shadowing dan resolusi nama;
  \item menerapkan aturan tipe pada ekspresi dan pernyataan;
  \item menangani konversi tipe dan pelaporan kesalahan;
  \item menyelesaikan M4: symbol table dan type checker.
\end{enumerate}
\end{outcomebox}

\section*{Materi Pokok}
\begin{itemize}[leftmargin=*]
  \item Struktur dan operasi tabel simbol.
  \item Implementasi hash table dan chaining.
  \item Nested scope, shadowing, dan resolusi nama.
  \item Sistem tipe dan pemeriksaan tipe.
  \item Konversi tipe dan pelaporan kesalahan.
\end{itemize}

\subfile{section-01}
\subfile{section-02}
\subfile{section-03}
\subfile{section-04}
\subfile{section-05}

\section*{Contoh}
\addcontentsline{toc}{section}{Contoh}

\begin{contoh}[Lookup pada scope bersarang]
Pada program
\begin{lstlisting}[language=C]
int x = 1;
{
    int y = 2;
    print(x + y);
}
\end{lstlisting}
saat memeriksa \code{x + y}, tabel mencari \code{x} di scope blok
(tidak ada), lalu di scope global (ditemukan, tipe \code{int}).
\code{y} ditemukan di scope blok (tipe \code{int}). Hasil
\code{x + y} bertipe \code{int}.
\end{contoh}

\begin{contoh}[Shadowing]
\begin{lstlisting}[language=C]
int x = 1;
{
    int x = 2;
    print(x);
}
print(x);
\end{lstlisting}
\code{print(x)} di dalam blok mencetak 2; \code{print(x)} di luar
blok mencetak 1. Tabel simbol menyimpan dua simbol \code{x} pada
scope berbeda.
\end{contoh}

\begin{contoh}[Pemeriksaan tipe ekspresi]
Untuk \code{int a; float b; a + b}, \code{check\_expr} pada
identifier \code{a} mengembalikan \code{int}, \code{b}
mengembalikan \code{float}. Aturan \code{type\_arith(int, float)}
mengembalikan \code{float}.
\end{contoh}

\begin{contoh}[Deklarasi ganda]
\begin{lstlisting}[language=C]
int main() {
    int x = 1;
    int x = 2;
    return 0;
}
\end{lstlisting}
Saat deklarasi kedua diproses, \code{lookup\_current} menemukan
\code{x} yang sudah ada. Pemeriksa melaporkan
\code{identifier 'x' sudah dideklarasikan} dan tidak menambahkan
simbol baru.
\end{contoh}

\begin{contoh}[Kesalahan konversi]
\begin{lstlisting}[language=C]
int main() {
    int x = 3.5;
    return 0;
}
\end{lstlisting}
Tipe inisialisasi adalah \code{float}, sedangkan tipe deklarasi
adalah \code{int}. Karena \code{float} tidak assignable ke
\code{int} tanpa cast, pemeriksa melaporkan
\code{tidak dapat menetapkan float ke int}.
\end{contoh}

\section*{Aktivitas Pembelajaran dan Praktikum}
\addcontentsline{toc}{section}{Aktivitas Pembelajaran dan Praktikum}

\begin{table}[H]
\centering
\caption{Rencana aktivitas Bab XV}
\label{tab:ch15-aktivitas}
\begin{tabularx}{\textwidth}{@{}llX@{}}
\toprule
\textbf{Sesi} & \textbf{Bentuk} & \textbf{Kegiatan} \\
\midrule
1 & Ceramah & Tabel simbol dan scope \\
2 & Live coding & Implementasi \code{symtab.c} \\
3 & Latihan & Pemeriksaan tipe ekspresi \\
4 & Praktikum & M4: symbol table dan type checker \\
\bottomrule
\end{tabularx}
\end{table}

\subsection*{M4 / Praktikum 4: Symbol Table dan Type Checker}
\textbf{Tujuan.} Mengimplementasikan tabel simbol bersarang dan
pemeriksa tipe untuk Subset-C.

\textbf{Keluaran.}
\begin{enumerate}[label=\arabic*.,leftmargin=*]
  \item \file{symtab.h} dan \file{symtab.c}.
  \item \file{types.h}, \file{types.c}, dan \file{typecheck.c}.
  \item Driver \file{main\_semantic.c}.
  \item Minimal 20 kasus uji semantik.
  \item Keluaran diagnosis untuk program tidak valid.
  \item Laporan singkat.
\end{enumerate}

\textbf{Langkah kerja.}
\begin{enumerate}[label=\arabic*.,leftmargin=*]
  \item Implementasikan \code{define}, \code{lookup}, dan scope.
  \item Implementasikan aturan tipe.
  \item Traversal AST untuk deklarasi dan ekspresi.
  \item Tangani \code{TYPE\_ERROR}.
  \item Uji program valid dan tidak valid.
\end{enumerate}

\textbf{Rubrik.}
\begin{table}[H]
\centering
\caption{Rubrik Praktikum 4}
\label{tab:ch15-rubrik}
\begin{tabularx}{\textwidth}{@{}lXr@{}}
\toprule
\textbf{Aspek} & \textbf{Indikator} & \textbf{Bobot} \\
\midrule
Kebenaran & Scope, lookup, dan tipe benar & 40\% \\
Kualitas kode & Struktur, penamaan, komentar & 30\% \\
Kelengkapan & Deklarasi, ekspresi, fungsi & 20\% \\
Pengujian & Kasus normal, shadowing, error & 10\% \\
\bottomrule
\end{tabularx}
\end{table}

\textbf{Definisi selesai.}
\begin{checklist}
  \item Tabel simbol mendukung scope bersarang.
  \item Pemeriksa tipe menangani promosi dan error.
  \item 20 kasus uji tersedia.
  \item Diagnosis menyebut baris, nama, dan jenis pelanggaran.
  \item Laporan diserahkan.
\end{checklist}

\section*{Latihan dan Refleksi}
\addcontentsline{toc}{section}{Latihan dan Refleksi}

\subsection*{Latihan}
\begin{enumerate}[label=\arabic*.,leftmargin=*]
  \item Sebutkan atribut yang disimpan tabel simbol.
  \item Jelaskan perbedaan \code{lookup} dan \code{lookup\_current}.
  \item Mengapa scope perlu didorong dan dilepas?
  \item Jelaskan shadowing dengan contoh.
  \item Rancang hash function sederhana untuk identifier.
  \item Tentukan tipe \code{int + char * float}.
  \item Jelaskan aturan konversi implisit.
  \item Rancang pesan kesalahan untuk \code{int x = 3.5;}.
\end{enumerate}

\subsection*{Refleksi}
\begin{reflectionbox}
\begin{enumerate}[label=\arabic*.,leftmargin=*]
  \item Bagian mana dari tabel simbol yang paling menantang?
  \item Apakah Anda dapat menjelaskan mengapa \code{lookup} dimulai
        dari scope dalam?
  \item Bagaimana Anda memastikan pesan kesalahan tidak berulang?
\end{enumerate}
\end{reflectionbox}

\section*{Asesmen}
\addcontentsline{toc}{section}{Asesmen}

Bab XV memuat Praktikum 4 dengan bobot 12\% sesuai RPS. Asesmen
mencakup tabel simbol, scope, dan pemeriksaan tipe.

\begin{table}[H]
\centering
\caption{Instrumen asesmen Bab XV}
\label{tab:ch15-asesmen}
\begin{tabularx}{\textwidth}{@{}lXl@{}}
\toprule
\textbf{Instrumen} & \textbf{Indikator} & \textbf{Bobot} \\
\midrule
Praktikum 4 & Symbol table, type checker, dan keluaran diagnosis &
12\% \\
Kuis & Scope, shadowing, dan aturan tipe & formatif \\
Review kode & Kualitas implementasi & bagian P4 \\
\bottomrule
\end{tabularx}
\end{table}

Kriteria kelulusan: tabel simbol mendukung scope bersarang,
pemeriksa tipe menangani promosi dan error, dan 20 kasus uji lulus.

\section*{Checklist Pencapaian Kompetensi}
\addcontentsline{toc}{section}{Checklist Pencapaian Kompetensi}

\begin{checklist}
  \item Saya dapat menjelaskan struktur tabel simbol.
  \item Saya dapat mengimplementasikan \code{define} dan
        \code{lookup}.
  \item Saya dapat mengelola scope bersarang.
  \item Saya dapat menjelaskan shadowing.
  \item Saya dapat menerapkan aturan tipe.
  \item Saya dapat menangani konversi implisit.
  \item Saya dapat menghasilkan pesan kesalahan tipe.
  \item Saya dapat menghasilkan berkas kode dan keluaran diagnosis.
  \item Saya sudah menyelesaikan M4/Praktikum 4.
\end{checklist}

\section*{Rangkuman}
\addcontentsline{toc}{section}{Rangkuman}

Pada bab ini telah dibahas:
\begin{itemize}[leftmargin=*]
  \item struktur dan operasi tabel simbol;
  \item implementasi hash table dan chaining;
  \item nested scope, shadowing, dan resolusi nama;
  \item sistem tipe dan pemeriksaan tipe;
  \item konversi tipe, pelaporan kesalahan, dan keluaran diagnosis;
  \item berkas kode symbol table dan type checker.
\end{itemize}

Bab XVI akan membahas AST, validasi kode, dan penanganan error
secara menyeluruh. Kita akan menggabungkan parser, tabel simbol, dan
pemeriksa tipe menjadi satu validator yang memberikan diagnostik
informatif.

\section*{Evaluasi}
\addcontentsline{toc}{section}{Evaluasi}

\begin{enumerate}[label=\arabic*.,leftmargin=*]
  \item Sebutkan lima atribut simbol.
  \item Jelaskan perbedaan \code{lookup} dan
        \code{lookup\_current}.
  \item Jelaskan shadowing dengan contoh.
  \item Rancang hash table untuk 100 identifier.
  \item Tentukan tipe \code{int + char * float}.
  \item Jelaskan aturan konversi implisit.
  \item Rancang pesan kesalahan untuk deklarasi ganda.
  \item Mengapa propagasi error penting?
\end{enumerate}

\subsection*{Kunci dan Rubrik Singkat}
\begin{itemize}[leftmargin=*]
  \item Soal 1--3: ketepatan konsep (70\%), contoh (30\%).
  \item Soal 4--6: rancangan dan aturan (70\%), penjelasan (30\%).
  \item Soal 7--8: kualitas pesan dan argumentasi (60\%),
        kelengkapan (40\%).
\end{itemize}


\ifSubfilesClassLoaded{%
  \bibliographystyle{apalike}%
  \bibliography{../../references}%
}{}

\end{document}
